\id{МРНТИ 50.47.02}{}

\begin{header}
\swa{Б.Б.~Оразбаев, Ж.А.~Кужуханова, Г.А.~Ускенбаева, Б.А.~Серимбетов, Л.Т.~Курмангазиев}{РАЗРАБОТКА МОДЕЛЕЙ ПРОЦЕССОВ ПЕРВИЧНОЙ ПЕРЕРАБОТКИ НЕФТИ В АТМОСФЕРНОМ БЛОКЕ ПРИ ДЕФИЦИТЕ И НЕЧЕТКОСТИ ИСХОДНОЙ ИНФОРМАЦИИ}

\tsp{1}Б.Б.~Оразбаев,
\tsp{2}Ж.А.~Кужуханова\envelope,
\tsp{1}Г.А.~Ускенбаева,
\tsp{2}Б.А.~Серимбетов,
\tsp{3}Л.Т.~Курмангазиев
\end{header}

\begin{affil}
\tsp{1}Евразийский национальный университет им. Л.Н. Гумилева, Астана, Казахстан,

\tsp{2}Казахский университет технологии и бизнеса им.К.Кулажанова, Астана, Казахстан,

\tsp{3}Атырауский университет им. Х. Досмухамедова, Атырау, Казахстан

\corrauthor{Корреспондент-автор:kuzhuhanova-zhad@mail.ru}
\end{affil}

Данная работа посвящена исследованию и решению задач моделирования
сложных систем на примере установки атмосферной перегонки нефти (УАВ) в
условиях ограниченности и неопределенности доступной информации.
Поскольку многие реальные технологические системы нефтеперерабатывающей,
нефтехимической и других отраслей промышленности характеризуются
неопределенностью и размытостью информации, необходимой для разработки
моделей, решение вопросов дефицита и нечеткости информации является
актуальной научной и практической задачей. Актуальность и важность
решения данной проблемы для промышленной практики также обусловлены
необходимостью повышения производительности и эффективности за счет
оптимизации режимов работы производственных систем на основе их моделей.
На практике при разработке моделей многих производственных объектов
часто возникают проблемы, связанные с нечеткостью части исходной
информации, влияющей на важные показатели работы объекта. Целью
исследования является разработка системного метода синтеза
статистических нечетких моделей сложных объектов в условиях дефицита и
нечеткости исходной информации. Далее с использованием предложенного
метода разработаны различные модели атмосферного блока установки
первичной переработки нефти. В данном случае статистические модели
разработаны на основе традиционных методов. При четких входных, режимных
параметрах и нечетких выходных параметрах на основе предложенного метода
разработаны нечеткие модели атмосферного блока, определяющие качество
выпускаемой продукции. А при нечетких входных, режимных и выходных
параметрах объекта качество целевого продукта атмосферного блока
определяется на основе методов экспертной оценки.

{\bfseries Ключевые слова:} установка первичной переработки нефти,
атмосферно-вакуумный блок, нечеткая информация, бензин, колонна.

\begin{header}
АНЫҚ ЕМЕС АҚПАРАТТЫҢ ЖЕТІСПЕУШІЛІГІ КЕЗІНДЕ АТМОСФЕРАЛЫҚ БЛОКТА МҰНАЙДЫ БАСТАПҚЫ ӨҢДЕУ ПРОЦЕСТЕРІНІҢ МОДЕЛДЕРІН ӘЗІРЛЕУ

\tsp{1}Б.Б.~Оразбаев,
\tsp{2}Ж.А.~Кужуханова\envelope,
\tsp{1}Г.А.~Ускенбаева,
\tsp{2}Б.А.~Серимбетов,
\tsp{3}Л.Т.~Курмангазиева
\end{header}

\begin{affil}
\tsp{1}Л.Н.Гумилев атындағы Еуразия Ұлттық университеті, Астана, Казахстан,

\tsp{2}Қ. Құлажанов атындағы Қазақ бизнес және технология университеті, Астана, Казахстан,

\tsp{3}Х. Досмухамедов атындағы Атырау университеті, Атырау, Казахстан,

e-mail: kuzhuhanova-zhad@mail.ru
\end{affil}

Бұл жұмыс қолжетімді ақпараттың шектеулілігі мен анық еместігі
жағдайында мұнайды бастапқы өңдеу қондырғысының атмосфералық айдау блогы
(ААБ) мысалында күрделі жүйелерді модельдеу мәселелерін зерттеуге және
шешуге арналған. Себебі мұнай өңдеу, мұнай-химия және басқа да өнеркәсіп
салаларындағы көптеген нақты технологиялық жүйелер модельдерді әзірлеу
үшін қажет ақпараттың анық еместігімен және бұлыңғырлығымен сипатталады.
Ақпарат тапшылығы мен бұлыңғырлығы мәселелерін шешу -- өзекті ғылыми әрі
практикалық міндет болып табылады.Бұл мәселені өнеркәсіптік практикада
шешудің өзектілігі мен маңыздылығы өндірістік жүйелердің жұмыс
режимдерін олардың модельдері негізінде оңтайландыру арқылы өнімділікті
және тиімділікті арттыру қажеттілігімен де айқындалады. Тәжірибеде
көптеген өндірістік объектілердің модельдерін әзірлеу кезінде бастапқы
ақпараттың бір бөлігінің бұлыңғыр сипатына байланысты қиындықтар жиі
туындайды, ол объектінің жұмыс көрсеткіштеріне әсер етеді.Зерттеудің
мақсаты -- бастапқы ақпарат тапшылығы мен бұлыңғырлығы жағдайында
күрделі объектілердің статистикалық және бұлыңғыр модельдерін
синтездеудің жүйелік әдісін әзірлеу. Ұсынылған әдіс негізінде мұнайды
бастапқы өңдеу қондырғысының атмосфералық блогының әртүрлі модельдері
жасалды. Бұл жағдайда статистикалық модельдер дәстүрлі әдістер негізінде
құрылды.Кіріс және режимдік параметрлер анық, ал шығыс параметрлері
бұлыңғыр болған жағдайда, ұсынылған әдіс негізінде өнім сапасын
анықтайтын атмосфералық блоктың бұлыңғыр модельдері жасалады. Ал егер
объектінің кіріс, режимдік және шығыс параметрлері бұлыңғыр болса,
атмосфералық блоктың мақсатты өнімі сапасы эксперттік бағалау әдістері
арқылы анықталады.

{\bfseries Түйін сөздер:} мұнайды бастапқы өңдеу қондырғысы, атмосфералық
вакуумдық қондырғы, анық емес ақпарат, бензин, колонна.

\begin{header}
DEVELOPMENT OF MODELS OF OIL PRIMARY REFINING PROCESSES IN AN ATMOSPHERIC UNIT IN THE CURRENT CONDITION OF UNCERTAIN INFORMATION

\tsp{1}B.B.~Orazbayev,
\tsp{2}Zh.A.~Kuzhukhanova\envelope,
\tsp{1}G.A.~Uskenbayeva,
\tsp{2}B.A.~Serimbetov,
\tsp{3}L.Т.~Kurmangaziyeva
\end{header}

\begin{affil}
\tsp{1}L.N. Gumilyov Eurasian National University, Astana, Kazakhstan,

\tsp{2} K.Kulazhanov Kazakh University of Technology and Business, Astana, Kazakhstan,

\tsp{3}Kh.Dosmukhamedov Atyrau University, Atyrau, Kazakhstan,

e-mail: kuzhuhanova-zhad@mail.ru
\end{affil}

This work is devoted to the study and solution of problems related to
modeling complex systems, using the atmospheric distillation unit of a
crude oil processing plant as an example, under conditions of limited
and uncertain available information. Since many real technological
systems in oil refining, petrochemical, and other industrial sectors are
characterized by uncertainty and vagueness of the information required
for model development, addressing the issues of information deficiency
and fuzziness is a relevant scientific and practical challenge. The
relevance and importance of solving this problem for industrial practice
also lie in the need to increase productivity and efficiency through the
optimization of operating modes of production systems based on their
models. In practice, when developing models of many production
facilities, problems often arise associated with the fuzzy nature of
some of the initial information, affecting important performance
indicators of the facility. The purpose of the study is to develop a
system method for synthesizing statistical, fuzzy models of complex
objects under conditions of deficiency and fuzziness of initial
information. Then, using the proposed method, various models of the
atmospheric block of the primary oil refining unit were developed. In
this case, statistical models were developed based on traditional
methods. With crisp input, mode parameters and fuzzy output parameters,
fuzzy models of the atmospheric block are developed based on the
proposed method, determining the quality of the manufac\-tured products.
And with fuzzy input, mode and output parameters of the object, the
quality of the target product of the atmospheric block is determined
based on expert assessment methods.

{\bfseries Keywords:} primary oil refining unit, atmospheric vacuum unit,
fuzzy information, gasoline, column.

\begin{multicols}{2}
{\bfseries Введение.} Процесс первичной переработки нефти, в частности в
атмосферном блоке, является неотъемлемой частью нефтехимической
промышленности, обеспечивающей разделение нефти на фракции, которые
используются для дальнейшей переработки и получения различных
нефтехимических продуктов. Атмосферная перегонка, основная операция в
атмосферном блоке, представляет собой процесс разделения нефти на
фракции с помощью разницы в температурах кипения компонентов.

Современные технологии переработки нефти требуют точного и эффективного
управления процессами в атмосферном блоке для обеспечения высокой
производительности и качества продуктов. В этом контексте важным
аспектом является создание математических моделей, способных описывать
процессы переработки нефти с учетом всех факторов, влияющих на их
протекание. Такие модели должны не только обеспечивать точность
прогнозов, но и быть адаптированными к реальным условиям производства,
где часто возникают неопределенности, вариации в составе нефти и внешние
воздействия {[}1{]}.

Технологические процессы первичной переработки нефти, как и другие
процессы нефтепереработки, нефтехимии, металлургии и других производств,
протекают в сложных химико-технологических системах, состоящих из разных
блоков. Например, процессы первичной переработки нефти, которые
протекают в технологическая схема атмосферного блока
ЭЛОУ-АТ-2Атырауского НПЗ, состоит из блоков ЭЛОУ, АТ и газореагентного
хозяйства {[}2{]}.

Блок ЭЛОУ предназначен для проведения процессов обезвоживания и
обессоливания сырой нефти на электродегидраторах. Блок атмосферной
перегонки АТ на основе нагревания, испарения и фракционирования нефти
обессоленной и обезвоженной нефти из блока ЭЛОУ, разделяет ее на
отдельные фракции, а также проводит конденсации паров дистиллятов. Блок
ГРХ предназначен для сбора, компаундирование и упорядоченная раздача
топливных газов, для слива, хранения, приготовления растворов редкого
натра необходимых концентраций и раздачи приготовленных растворов на
технологические установки завода.

Выбор режима работы блока АТ в основном зависит от фракционного состава
нефти, от содержания в ней бензиновых фракций и газов, а также от
значения входных, режимных параметров основных агрегатов данного блока.
К основным агрегатам блока АТ можно отнести ректификационную колонну
отбензинивания нефти К-1 и атмосферную колонну К-2.

Для создания систему управления процессом первичной переработки и выбора
эффективного режима работы блока АТ необходимо разработать
взаимосвязанных математических моделей основных колонн блока, т.е.
ректификационной колонны К-1иатмосферной колонныК-2. На практике
разработка математических моделей этих колонн усложняется из-за дефицита
и нечеткости необходимой исходной информации для создания моделей
{[}3{]}. В производственных условиях определить фракционного состава
перерабатываемой нефти, значения некоторых важных входных, режимных
параметров указанных колонн, влияющие на их выходные параметры
измерительными приборами очень сложно или невозможно.

На практике такими количественно трудно описываемыми объектами часто
хорошо управляет эксперты, управляющие именно основе своего опыта,
знания и интуиции. Их опыт, знания и интуиция, выраженная на
естественном языке как нечеткая информация, на основе методов экспертных
оценок, теорий нечетких множеств можно собрать, формализовать и
использовать для описания режимов работы и разработки моделей объектов.
В этой связи в настоящее время вопросы разработки адекватных моделей
сложных, количественной трудно описываемых объектов, таких как К-1, К-2,
характеризующие дефицитом и нечеткостью исходной информации является
актуальной научно-практической проблемной. В данном исследовании для
эффективного решения данной проблемы разработки моделей сложных объектов
в условиях дефицита и нечеткости исходной информации проблем
предлагается подход, основанные на использование доступной нечеткой и
другой информации различного характера.

Обзор и анализ работ по теме исследования, посвященных к вопросам
разработки моделей сложных, производственных объектов и технологических
процессов. В исследовательских работах предложены различные подходы к
разработке моделей и моделированию технологических объектов в различных
условиях в т. ч. с использованием нечеткой информации. В частности, в
исследованиях под руководством Кафарова В. {[}4{]} разработаны
детерминированные модели основных процессов различных производств,
протекающие в технологических агрегатах на основе теоретических
сведений, кинетики процессов. В качестве преимуществ детерминированных
моделей, предложенных в этих исследованиях типовых агрегатов и процессов
различных производств, можно отметить их универсальность и теоретический
обоснованность, так как они разработаны на основе фундаментальных
законно в природы.

Недостатки этих исследований в том, что такие детерминированные модели
для сложных технологических объектов, как исследуемые колонны К-1,
К-2невозможно или очень сложна, разработка их становится экономически
нецелесообразной. Это связана тем, что такие сложные объекты
характеризуются дефицитом, и/или отсутсвие необходимых теоретических
сведений для разработки детерминированных моделкей, они характеризуются
нечеткостью исходной информации. При этом принятия различных допущений и
идеализация условии работы сложных объектов для разработки
детерминированных моделей при дефиците, нечеткости исходной информации
приводят к получению не адекватных, не пригодных моделей для применения
на практике.

Авторы работ {[}5{]} исследовали и предложили подходы к разработке
моделей и управления на их основе объектами в условиях нечеткости
некоторой части исходной информации, основанные на применение методов
экспертной оценки и теорий нечетких множеств. Предложенные подходы к
разработке моделей на основе множества уровня (α) позволяют
разрабатывать нечеткие модели при четких значениях входных, режимных
параметров и нечетких значениях выходных параметров объекта. Но в этих и
других исследованиях разработки моделей на основе нечеткой информации не
исследованы и не решены проблемы разработки моделей в условиях нечетких
входных, режимных, и выходных параметров объекта. Кроме того, в
анализированных и в других исследованиях по разработке моделей не
исследованы проблемы разработки взаимосвязанных систем моделей
технологических агрегатов, как колонны К-1 и К-2 атмосферному блоку,
связанные через печи подогрева. Такая ситуация мотивировала данного
исследования. Нами для решения проблем разработки систему моделей
нечетко описываемых объектов применяется другой подход, основанные
методы системного анализа, экспертных оценок и теорий нечетких множеств.

{\bfseries Материалы и методы.} Объектом исследования является колонны К-1
и К-2 атмосферному блоку установки первичной переработки нефти ЭЛОУ-АВТ
Атырауского НПЗ, которая характеризуется нечеткостью исходной
информации. Материалами исследования являются технологические схемы,
описание и технологический регламент работы атмосферного блока установки
первичной переработки Атырауский НПЗ, а также статистические данные и
нечеткая информация о состояние и его режимах его работы. На рисунке 1
приведена технологическая схема атмосферного блока, выделенная из общей
схемы установки ЭЛОУ-АВТ2 Атырауского НПЗ.

Приведем описания процесса первичной переработки нефти в соответствие с
приведенной технологической схемой. Сырая нефть подогревается и подается
наблок обессоливания и электрообезвоживания (блок ЭЛОУ). Затем
обессоленная нефть нагревается и подается в предварительный испаритель
блока атмосферной перегонки, где частично отбензинивается.
Нижний продукт (частично отбензиненная нефть) после нагрева в печи
поступает в основную атмосферную колонну, где отбираются три боковые
погона - керосиновый (фракция 180-220\tsp{0}С), легкий
дизельный (фракция 220-280\tsp{0}С), тяжелый дизельный
(фракция 280-350\tsp{0}С). Нижний продукт (мазут), при
неработающем вакуумном блоке, после охлаждения поступает в
товарно-сырьевую базу для приготовления товарных мазутов {[}6{]}.

При включенном вакуумном блоке мазут после подогрева поступает в
вакуумную колонну и разделяется на следующие фракции: легкий вакуумный
газойль; тяжелый вакуумный газойль; затемненный продукт гудрон. Верхние
продукты колон К-1 и К-2 (бензиновая фракция НК-180\tsp{0}С)
смешиваются, отделяются от воды и частично от газов и подаются в
стабилизационную колонну. Верхний продукт этой колонны охлаждается и
выводится с установки в виде конденсата и газа. Нижний продукт -
стабильная бензиновая фракция НК-180\tsp{0}С поступает на
блок вторичной перегонки, где разделяется на узкие фракции
(фр.62\tsp{0}С, фр.62-105\tsp{0}С,
фр.105−140\tsp{0}С, фр.140-180\tsp{0}С) {[}7{]}.
\end{multicols}

\fig[0.7\textwidth]{i2/image15}[Рис.1 - Технологическая схема атмосферного блока ЭЛОУ-АТ-2]

\begin{multicols}{2}
Системный метод разработки моделей является ключевым подходом к
моделированию и анализу сложных объектов, особенно в условиях дефицита и
неопределённости исходной информации. Этот метод позволяет создавать
обоснованные и эффективные модели, которые помогают принимать решения в
ситуациях, когда стандартные подходы неприменимы из-за неопределённости,
нечёткости данных. Важно отметить, что такие модели могут включать
статистические и нечеткие элементы, каждый из которых играет свою роль и
имеет свои особенности в контексте обработки информации.

Разработка моделей для оценки объема бензина на выходе из атмосферного
блока с учетом входных и режимных параметров на основе
экспериментально-статистического подхода. Для разработки моделей
взаимосвязанных колонн К-1 и К-2 используются различные методы, в том
числе статистические подходы {[}8{]}, методы системного анализа {[}9{]},
экспертные оценки и теория нечетких множеств {[}10{]}.

Предложенный метод включает в себя несколько основных этапов.

1) На основе системного анализа собрать и обработать доступную
информацию статистического \(x_{i},\ i = \overline{1,n_{1}}\)и нечеткого
\({\widetilde{x}}_{i},\ i = \overline{n_{1} + 1,n}\)характера о режимах
работы взаимосвязанных объектов технологического комплекса, влияющие на
выходные параметры, которые характеризуют качество работы объекта.
Выходные параметры объекта также могут быть четкими
\(y_{j},\ j = \overline{1,m_{1}}\) и нечеткими
\({\widetilde{y}}_{j},\ j = \overline{m_{1} + 1,m},\) оцениваемые
экспертами.

2) Генерировать критериев \(f_{k},\ k = \overline{1,K}\)выбора
эффективного типа моделей для каждого объекта технологического
комплекса.

3) Для каждого объекта технологического комплекса на основе характера
доступной исходной информации и критериев выбора
\(f_{k},\ k = \overline{1,K}\) методами экспертной оценки определить
эффективный тип разрабатываемой модели.

4) Синтезируются нечеткие модели для объекта с четкими входными и
нечеткими выходными параметрами. Структуру нечетких моделей на основе
идеи метода последовательного включения регрессоров можно
идентифицировать в виде:
\end{multicols}
\vspace{-0.5em}
\begin{equation}
\tilde{y}_{j} = \tilde{a}_{0j} + \sum_{i=1}^{n_{1}} \tilde{a}_{ij} x_{ij} + \sum_{i=1}^{n_{1}} \sum_{k=i}^{n_{1}} \tilde{a}_{ikj} x_{ij} x_{kj} + \dots, \quad i = \overline{1, n_{1}}, \ j = \overline{m_{1}+1, m},
\end{equation}

\begin{multicols}{2}
где \({\widetilde{a}}_{0j},{\widetilde{a}}_{ikj}\) - неизвестные нечеткие
параметры (регессионные коэффициенты), которых необходимо
идентифицировать. Первая сумма состьавляет линейную часть нечетких
моделей, последующие суммы - нелинейную часть.

5) Если модели удовлетворяют условиям адекватности, они могут быть
рекомендованы для практического использования при оптимизации и
управлении режимами работы объекта. В противном случае выявляются
причины их неадекватности, и процесс возвращается на соответствующие
этапы метода для устранения этих недостатков и обеспечения соответствия
моделей реальным условиям.

6) Идентифицируются нечеткие параметры моделей (1). Предлагается
следующий подход к идентификации нечетких параметров. Нечеткие модели на
основе множества уровня α преобразуются в набор четких моделей:
\end{multicols}
\vspace{-0.5em}
\begin{equation}
y_{j}^{\alpha_{l}} = a_{0j}^{\alpha_{l}} + \sum_{i = 1}^{n_{1}}{{a_{ij}^{\alpha_{l}}x}_{ij} + \sum_{i = 1}^{n_{1}}{\sum_{k = i}^{n_{1}}{{a_{ikj}^{\alpha_{l}}x}_{ij}x_{kj}}}},\ l = \overline{1,L},
\end{equation}

\begin{multicols}{2}
где \(\alpha_{l}\mathbf{,\ } = \overline{1,L}\) - множество
уровня(\emph{α)}.

Затем неизвестные значения параметров
\(a_{0j}^{\alpha_{l}},a_{ij}^{\alpha_{l}},\ a_{ikj}^{\alpha_{l}}\)полученных
четких моделей набор можно идентифицировать с помощью известных методов
параметрической идентификации. После идентификации параметров модели (2)
на α уровнях, например на основе метод наименьших квадратов, полученных
значений параметров необходимо объединить по формуле {[}11{]}:
\end{multicols}
\vspace{-0.5em}
\begin{equation}
\begin{aligned}
{\widetilde{a}}_{ij} = \bigcup_{\alpha \in [0.5\div 1]} a_{ij}^{\alpha_{l}} \quad &\text{or} \quad \mu_{{\widetilde{a}}_{ij}}(a_{ij}) = \sup_{\alpha \in [0.5, 1]} \min \{ \alpha_{l}, \mu_{a_{ij}}^{\alpha_{l}}(a_{ij}) \}, \\
\text{где} \quad &a_{ij}^{\alpha_{l}} = \left\{ a_{ij} \mid \mu_{{\widetilde{a}}_{ij}}(a_{ij}) \geq \alpha \right\}
\end{aligned}
\end{equation}

\begin{multicols}{2}
7) Проверяется адекватность разработанных моделей. При этом в качестве
критериев адекватности можно выбрать:
\vspace{-0.5em}
\[R = \min{\sum_{j = 1}^{m}\left( y_{j}^{M} - y_{j}^{E} \right)^{2}} \leq R_{D},\]

где \(y_{j}^{M}\)и \(y_{j}^{E}\)-соответственно модельные, рассчитанные
по моделей и реальные, полученные, например экспериментальным путем на
объекте, значения выходным параметров; \(R_{D}\)-допустимое отклонение
между \(y_{j}^{M}\)и \(y_{j}^{R}.\)

8) Если модели соответствуют условиям адекватности, их можно
рекомендовать для практического применения в процессе оптимизации и
управления режимами работы моделируемого объекта. В случае выявления
неадекватности моделей проводится анализ причин этого несоответствия,
после чего осуществляется возврат к соответствующим этапам метода для
устранения проблем и обеспечения адекватности разрабатываемых моделей.

{\bfseries Обсуждение и результаты.} На основе системного метода разработки
статистических, нечетких моделей сложных объектов в условиях дефицита и
неопределенности исходной информации, предложенного в предыдущем
разделе, синтезированы модели колонн К-1 и К-2 для атмосферного блока
установки первичной переработки нефти. В соответствии с пунктом 4
предложенного метода была определена структура статистической модели
колонны К-1, которая описывает выход бензина фракции
180\tsp{0}С и отбензиненной нефти, подаваемой в колонну К-2
\end{multicols}
\begin{equation}
y_{j} = a_{0j} + \sum_{i=1}^{3} a_{ij}x_{i} + \sum_{i=1}^{3} \sum_{k=i}^{3} a_{ikj}x_{i}x_{k}, \quad j=1,2
\end{equation}

где \(y_{1}\) -объем бензина фракции 180 $^\circ$С, подаваемый
на стабилизационной колонны,\(y_{2}\) - объем отбензиновой нефти,
подаваемой в К-2; \(x_{i},\ i = \overline{1,3}\) входные, режимные
параметры К-1, \(x_{1} -\) объем сырья (обессоленной
нефти), \(x_{2} -\) температура в К-1, \(x_{3} -\) давление в К-1.

После определения неизвестных параметров моделей (4) с использованием
экспериментально-статистических данных и пакета программ REGRESS, в
котором реализован метод наименьших квадратов, были получены модели для
расчета объема выхода продуктов из колонны К-1:

\begin{equation}
\begin{aligned}
y_{1} &= 11.1 + 0.008615385x_{1} + 0.029473684x_{2} - 1.555555556x_{3} + 0.000026509x_{1}^{2} \\
      &\quad + 0.000077562x_{2}^{2} + 0.864197531x_{3}^{2}
\end{aligned}
\end{equation}

\begin{equation}
\begin{aligned}
y_{2} &= 118.7 + 0.091384615x_{1} + 0.312631579x_{2} - 16.500000000x_{3} + 0.000281183x_{1}^{2} \\
      &\quad + 0.000822715x_{2}^{2} + 9.166666667x_{3}^{2}.
\end{aligned}
\end{equation}

\begin{multicols}{2}
В соответствии со структурой (1) предложенного системного метода
идентифицирована следующая структура нечеткой модели, оценивающей конец
кипения, бензина фракции «не выше»,
т.е.\(\widetilde{<}180\)$^\circ$С:
\end{multicols}

\[{\widetilde{y}}_{3} = {\widetilde{a}}_{0} + {\widetilde{a}}_{1}x_{2} + {\widetilde{a}}_{2}x_{3} + {\widetilde{a}}_{3}x_{2}^{2} + {\widetilde{a}}_{4}x_{3}^{2}.\ \ \ \ \ \ \ \ \ \ \ \ \ \ \ \ \ \ \ \ \ \ \ \ \ \ \ \ \ \ \ \ \ (7)\]

Продукт с низа колонны К-1, подогревается в печи П-2до необходимой
температуры (360\(-\)385\tsp{0}С) и подается в
ректификационную колонну К-2 для разделения на другие фракции бензина и
на мазут. Разработанные модели колонны К-2 на основе предложенного
системного метода, определяющие объемы бензиновых фракций
180--220\tsp{0},220−280\tsp{0} и
280−350\tsp{0}С, мазута фракции
<350\tsp{0}С имеют следующие структуры:

\[y_{j} = a_{0j} + \sum_{i = 4}^{6}{a_{ij}x_{i} +}\sum_{i = 4}^{6}{\sum_{k = i}^{6}{a_{ikj}x_{i}}x_{k},\ \ \ j = \overline{4,7,}\ \ \ \ \ \ \ \ \ \ \ \ \ \ \ \ \ \ \ \ \ \ \ \ \ \ \ \ \ \ (8)}\]

где\(y_{4} -\)объем бензиновых фракций180-220\tsp{0}С;
\(y_{5} -\)объем бензиновых
фракций220-280\tsp{0}С;\(y_{6} -\) объем керосино-газойлевых
фракций280-350\tsp{0}С; \(y_{7} -\)объем мазута фракции
350\tsp{0}С; \(x_{i},\ i = \overline{4,6} -\)входные,
режимные параметры К-2:\(x_{4} -\)объем сырья (частично отбензиновая
нефть из К-1); \(x_{5} -\)температура в К-2; \(x_{6} -\)давление в К-2.

После идентификации неизвестных параметров моделей (8) описанным выше
способом с помощью программы REGRESS, получены следующие модели
определения объема выхода продуктов из С-2 в зависимости от ее входных,
режимных параметров:

\[y_{4} = 11.2 + 0.00471381x_{4} + 0.02709677x_{5} - 1.4x_{6} + \ 0.00003174x_{4}^{2} + \ {0.000029136x}_{5}^{2}\ \ \  + \ 0.35x_{6}^{2},\ \ \ (9)\]

\[y_{5} = 36 + 0.016835017x_{4} + 0.0967741945 - 3.0x_{6} + 0.000113367x_{4}^{2} + {0.000104058x}_{5}^{2} + \ 1.25x_{6}^{2},\ \ (10)\]

\[y_{6} = 4 + \ 0.001683502x_{4} + 0.009677419x_{5} - 0.5x_{6} + 0.000011337x_{4}^{2}{+ 0.000010406x}_{5}^{2} + 0.125x_{6}^{\ 2},\ \ (11)\]

\[y_{7} = 48 + 0.025252525x_{4} + 0.145161290x_{5} + 1.5x_{6} + 0.000170051x_{4}^{2}{+ 0.000156087x}_{5}^{2} + 1.875x_{6}^{2}.\ \ \ \ \ \ \ \ (12)\]

Идентифицирована следующая структура нечетких моделей, позволяющий
оценить качества бензиновых фракций и мазута из К-2:

\[{\widetilde{y}}_{j} = {\widetilde{a}}_{0} + {\widetilde{a}}_{1}x_{5} + {\widetilde{a}}_{2}x_{6} + {\widetilde{a}}_{3}x_{5}^{2} + {\widetilde{a}}_{4}x_{5}x_{6} + {\widetilde{a}}_{5}x_{6}^{2}\widetilde{<}b_{j},\ j = \overline{8.11}\ \ \ \ \ \ \ \ \ \ \ \ (13)\]

где
\({\widetilde{y}}_{8},{\widetilde{y}}_{9},{\widetilde{y}}_{10}и{\widetilde{y}}_{11}\mathbf{-}\)соответственно,
бензиновые фракции с концами кипения\(b_{8}\widetilde{\leq}\)200,
\(b_{9}\widetilde{\leq}\)280,
\(b_{10}\widetilde{\leq}\)350\tsp{0}С и мазут фракции
началом
кипения\(b_{11}\widetilde{\geq}\)350\tsp{0}С.

Затем аналогично для нечеткой модели (7) формуле (2) представлены
набором четких моделей на множестве уровней \emph{α}= \{0.5;0.8;1\} и с
помощью программы REGRESS идентифицированы регрессионные коэффициенты
полученных четких моделей на (\emph{α)} уровнях. Параметры остальных
нечетких моделей, оценивающие качества других бензиновых фракций и
мазута идентифицированы аналогичным образом.

После объединения идентифицированных параметров на (\emph{α)} уровнях по
формуле (7) получены следующие модели для определения конца кипения
бензиновых фракции с колонны К-2:

\({\widetilde{y}}_{8} = 100.0 + \ 0.153846154x_{5} + 30.00x_{6} + 0.001183432x_{5}^{2}\widetilde{<}200,\ \ \ \ \ \ \ \ \ \ \ \ (14)\)

\[{\ \ \ \ \ \widetilde{y}}_{9} = \ 137.50 + 0.110x_{5} + 41.250x_{6} + 0.000440x_{5}^{2}\widetilde{<}280.\ \ \ \ \ \ \ \ \ \ \ \ \ \ \ \ \ \ \ \ \ \ \ \ \ \ \ \ \ \ \ \ \ \ \ \ (15)\]

\({\widetilde{y}}_{10} = \ 175 + 0.102941176x_{5} + 52.50x_{6} + 0.000302768x_{5}^{2}\widetilde{<}350,\ \ \ \ \ \ \ \ \ \ \ \ (16)\)

\[{\widetilde{\ y}}_{11} = \ 175.0 + \ 0.106060606x_{5} + 52.50x_{6} + 0.000321396x_{5}^{2}\widetilde{>}350.\ \ \ \ \ \ \ \ \ (17)\]

В полученных выше моделях колонн К-1, К-2 атмосферного блока установки
первичной переработки нефти (9) - (8), (14)-(12) и (14)-(17) регрессоры,
которые имеют нулевые коэффициенты, как не влиящие на выходые параметры
пренебрежены и в указанных моделях не приведены. Результаты
моделирования на основе разработанных и известных моделей для удобства
анализа и сравнения занесены в таблицу 1.

{\bfseries Таблица 1 - Результаты моделирования процесса первичной
переработки нефти с известными моделями {[}11{]} и разработанной
системой различных моделей с использованием дополнительной нечеткой
информации и экспериментальных реальных данных с блока АТ Атырауского
НПЗ}

%% \begin{longtable}[]{@{}
%%   >{\raggedright\arraybackslash}p{(\linewidth - 6\tabcolsep) * \real{0.5516}}
%%   >{\centering\arraybackslash}p{(\linewidth - 6\tabcolsep) * \real{0.1345}}
%%   >{\centering\arraybackslash}p{(\linewidth - 6\tabcolsep) * \real{0.1644}}
%%   >{\centering\arraybackslash}p{(\linewidth - 6\tabcolsep) * \real{0.1495}}@{}}
%% \toprule\noalign{}
%% \begin{minipage}[b]{\linewidth}\centering
%% {\bfseries Объемы и качественные показателей продуктов}
%% \end{minipage} & \begin{minipage}[b]{\linewidth}\centering
%% {\bfseries Известные модели {[}11{]}}
%% \end{minipage} & \begin{minipage}[b]{\linewidth}\centering
%% {\bfseries Разработанные системы моделей}
%% \end{minipage} & \begin{minipage}[b]{\linewidth}\centering
%% {\bfseries Реальные эксперимен-}
%% 
%% {\bfseries тальные данные}
%% \end{minipage} \\
%% \midrule\noalign{}
%% \endhead
%% \bottomrule\noalign{}
%% \endlastfoot
%% Объем бензина фракции 180\tsp{0}С с колонны С-1, \(y_{1}\);
%% & 27,2 & 28,7 & 28 \\
%% Объем от бензиновой нефти с колонны С-1\(y_{2}\); & 299 & 297 & 297 \\
%% Качество бензина фракции 180\tsp{0}С,
%% \({\widetilde{y}}_{3}\) & − & 175 & (178)\tsp{L} \\
%% Объем бензина фракции 180-220\tsp{0}С с колонны С-2,
%% \(y_{4};\) & 27,2 & 28,3 & 27.8 \\
%% Объем бензина фракции 220-280\tsp{0}С с колонны С-2,
%% \(y_{5};\) & 100 & 100 & 100 \\
%% Объем бензина фракции 280-350\tsp{0}С с колонны С-2,
%% \(y_{6};\) & 10.3 & 10 & 10 \\
%% Объем мазута с началом кипения 350\tsp{0}С с колонны С-2,
%% \(y_{7};\) & 151 & 149.5 & 150 \\
%% Качество бензина фракции 180-220\tsp{0}С, \(y_{8};\) & − &
%% 200,3 & (210)\tsp{L} \\
%% Качество бензина фракции 220-280\tsp{0}С, \(y_{9};\) & − &
%% 275 & (278)\tsp{L} \\
%% Качество бензина фракции 280-350\tsp{0}С, \(y_{10};\) & − &
%% 347 & (350)\tsp{L} \\
%% Качество мазута с началом кипения 350\tsp{0}С,
%% \({\widetilde{y}}_{11};\) & − & 358 & (360)\tsp{L} \\
%% Качество бензина с выхода блока АТ,\(y_{12};\) & − & 200 &
%% (200)\tsp{L} \\
%% Значения входных, режимных параметров блока АТ\(x_{1}^{*} -\)объем сырья
%% в К-1 (обессоленной нефти); & 325.7 & 325.7 & 325.7 \\
%% \(x_{2}^{*}\)--температура в К-1; & 193 & 190 & 190 \\
%% \(x_{3}^{*}\)--давление в К-1; & 1,85 & 1,80 & 1,80 \\
%% \(x_{4}^{*} -\)объем сырья в К-2; & 297 & 297 & 297 \\
%% \(x_{5}^{*}\)-- температура в К-2; & 312 & 310 & 311 \\
%% \(x_{6}^{*}\)-- давление в К-2; & 2,10 & 2 & 2 \\
%% \(x_{7}^{*}\)-- содержание хлористых солей в сырье блока АТ & − & 4 &
%% (4)\tsp{L} \\
%% \(x_{8}^{*}\)-- массовая доля серы в сырье блока АТ & − & 0.75 &
%% (0.75)\tsp{L} \\
%% \end{longtable}

\emph{Note:} - означает, что данным методом эти параметры не
определяются;
()\emph{\tsp{L}}--
эти параметры непосредственно не измеряются, оценивается в лаборатории с
участием человека.

{\bfseries Выводы.} Разработаны взаимосвязанные модели ректификационных
колонн К-1 и К-2 атмосферному блоку установки первичной переработки
нефти. Для определения объема нефтепродуктов на выходе атмосферного
блока с помощью предложенного системного метода были созданы
статистические нелинейные модели полиномиального типа (5) -(6), (9)
-(12). Структура и параметры этих моделей были идентифицированы с
использованием метода последовательного включения регрессоров и
экспериментально-статистических данных с применением программы REGRESS,
реализующей метод наименьших квадратов. Для оценки нечетких качественных
показателей бензиновых фракций и мазута с атмосферного блока были
разработаны нечеткие модели (8), (14) -(17). Структура этих моделей была
определена на основе модификации метода последовательного включения
регрессоров. Нечеткие параметры были идентифицированы путем
преобразования нечетких моделей в наборы четких моделей с использованием
множества уровня (\emph{α)} теории нечетких множеств, согласно формуле
(2). Затем с помощью экспериментально-статистических данных и программы
REGRESS были получены значения параметров для различных уровней α.
Объединив эти идентифицированные значения параметров по формуле (3),
были определены параметры моделей (8), (14) -(17).

Предлагаемый метод к разработке математических моделей основаан
комплексное использование доступной информации различного хаарактера, в
т. ч. нечеткой информации, предситавляющей собой знания, опыт и
соображения ЛПР, специалистов-экспертов по решаемой проблеме.
Предлагаемый метод отличается от существующих детерминированных,
статистияческих методов тем, что за счет использования доступной, в т.ч.
нечеткой информации позволяет разрабатывать эффективные и
высокоадекватные модели исследуемого объекта, функционирующих в условиях
неопределенностью.

За счет комплексного использования доступной информации различного
характера и гибридного подхода к разработке моделей сложного объекта
производства предлагаемый системный метод имеет эффект синергизма и
свойства эмердженности системы, которые является преимуществом
предлагаемого метода разработки сложных объектов. А эффект синергизма и
свойства эмердженности предлагаемого системного метода позволяет
разработать эффективные и высокоадекватные модели для системного
моделирования сложных производственных объектов, характеризующихся
дефицитом и нечеткостью исходеной информации.

Проведем анализ и сравнения результатов моделирования работы
атмосферного блока на основе известных моделей и разработанной системой
различных моделей с использованием дополнительной нечеткой информации,
приведенных в таблице 1, которые позволяет выделить следующие
преимущества предложенного метода моделирования:

1. Предложенный метод моделирования, основанный на разработанной системе
моделей с использованием нечеткой информации, позволяет увеличить объем
более важных целевых фракций бензина за счет применения опыта, знаний и
интуиции специалистов и экспертов. Как показано в результате
сравнительного анализа (таблица 1), объем более важных бензиновых
фракций 180\tsp{0}С и 180--220\tsp{0}С на выходе
атмосферного блока увеличивается на 1,5 и 1,1 т/час, что соответствует
росту на 5,51\% и 4,04\% соответственно, в то время как объемы менее
важных продуктов уменьшаются.

2. Предложенныйметод за счет использования нечеткую информацию,
представляющей собой опыт, знаня и интуцию ЛПРпозволяет определить
нечетко описываемых качественных показателейпроизводимых целевых
нефтепродуктов
(\({\widetilde{y}}_{3},\ {\widetilde{y}}_{8},{\widetilde{y}}_{9},{\widetilde{y}}_{10},{\widetilde{y}}_{11},{\widetilde{y}}_{12}\)),
которые не определяются известными методами. Как видно из результатов
сравнения (таблица 1).

3. При моделировании на основе разработанных моделей при меньших
значенияхтемпературыв и давления в К-1 и К-2
(\(x_{2},x_{3},x_{5\ }иx_{5})\)получены лучшие значения выходных целевых
продуктов. Таким образом разработанные модели позволяет получить больше
целевых продуктов при меньших энергозатрат, что обосновывает их
преимущества перед известными детерминированными моделями.

Рассмотрим результаты критического анализа и возможных ограничений
предлагаемого системного метода разработки моделей сложных
производственных объектов, которые могут возникнуть при применении его
на практике.

Предлагаемый системный метод разработки моделей сложных объектов в
условиях неопределенности предполагает наличие ЛПР и специалистов
экспертов по исследуемому объекту. В случае отсутствия таких ЛПР,
специалистов-экспертов метод может оказаться неэффективным или
непригодным для разработки эффективных моделей сложных, плохо
формализуемых объектов в условиях неопределенности. Но для таких
производственных объектов, как атмосферный блок установки первичной
переработки нефти Атырауского НПЗ, функционирующих долгие годы, как
правило, ЛПР и опытные специалисты, являющиеся экспертами по объекту
исследования, имеются.

Так как в предлагаемом системном методе проблемы неопределенности
решаются в основном за счет экспертной информации, для результатов
экспертной оценки практики ограничением к применению этого метода может
быть субъективизм экспертной оценки. Кроме того, часто необходимо
значительное время для обработки результатов экспертной оценки и
достижения согласованности мнений экспертов. Для решения этих проблем
предлагается модифицировать процедуры экспертизы, автоматизировать,
реализовав ее в компьютерной сети с помощью соответствующей программы.
При этом за счет автоматизации минимизируется время проведения
экспертизы обеспечения достижения согласованных мнений экспертов,
благодаря которой значительно снижается субъективизм экспертной оценки.

{\bfseries Литература}

1. Затонский А.В., Тугашова Л.Г. Управление атмосферной колонной малого
нефтеперерабатывающего завода с применением динамической
модели//Науковедение. - 2017. -Т.9(1). -С.1-13.

2. Кембаев А.Р. Табиғи мұнай және газды өңдеудегі жабдықтарды жөндеу//
Алматы: TechSmith.2023. -276c . ISBN 978-601-352-777-2.

3. Fayaz M., Ahmad S., Ullah I., Kim D. A Blended Risk Index Modeling
and Visualization Based on Hierarchical Fuzzy Logic for Water Supply
Pipelines Assessment and Management//Processes. -2018. -Vol.6(5):61. DOI
10.3390/pr6050061.

4. Кафаров В.В., Дорохов И.Н. Системный анализ процессов химической
технологии. основы стратегии: монография//Издательство Юрайт.
-2026.-499c. ISBN~978-5-534-06991-4.

5. Orazbayev B.B, Ospanov Ye.A., Orazbayeva K.N, Makhatova V.E.,
Urazgaliyeva M.K., Shagayeva A.B. Development of mathematical models of
R-1 reactor hydrotreatment unit using available information of various
types// \href{https://iopscience.iop.org/journal/1742-6596}{Journal of
Physics: Conference Series} 1399. -2019. -Vol.1399(4):044024. DOI
10.1088/1742-6596/1399/4/044024.

6. Гречухина А.А., Елпидинский А.А., Мингазов Р.Р., Плохова С.Е. Расчет
ректификационных колонн установок перегонки нефти. Учебное пособие
//Казанский национальный исследовательский технологический университет.
-2017.-92 c. ISBN\\
978-5-7882-2138-0.

7. Власов В.Г. Подготовка и переработка нефтей. -2-е изд.//
Инфра-Инженерия. -2026. -328 с. ISBN 978-5-9729-2808-8.

8. Orazbayev B., Ospanov Ye., Makhatova V., Salybek L., Abdugulova Z,
Kulmagambetova Z., Suleimenova S., Orazbayeva K. Methods of Fuzzy
Multi-Criteria Decision Making for Controlling the Operating Modes of
the Stabilization Column of the Primary Oil-Refining
Unit//Mathematics.-2023.-Vol.11(13): 2820. DOI
\href{https://doi.org/10.3390/math11132820}{10.3390/math11132820}.

9. Музипов Х.Н. Системы управления технологическими процессами добычи,
промысловой подготовки и транспорта нефти и газа. Учебное пособие для
вузов // Издательство Лань-Пресс. -2025. -268 c. ISBN 978-5-507-46261-2.

10. Anwar M.Z., Al-Kenani A.N., Bashir S., Shabir M. Pessimistic
Multigranulation Rough Set of Intuitionistic Fuzzy Sets Based on Soft
Relations//Mathematics. -2022. -Vol.10(5):685. DOI 10.3390/
math10050685.

11. Меркулов В.В. Химическая технология. Переработка нефти и газа:
Учебное пособие// Издательство Бастау. -2018. - 260 C. ISBN
978-601-7275-69-6.

{\bfseries References}

1. Zatonskij A.V., Tugashova L.G. Upravlenie atmosfernoj kolonnoj malogo
neftepererabatyvajushhego zavoda s primeneniem dinamicheskoj
modeli//Naukovedenie. - 2017. -T.9(1). -S.1-13. {[}in Russian{]}

2. Kembaev A.R. Tabiғi mұnaj zhәne gazdy өңdeudegі zhabdyқtardy
zhөndeu// Almaty: TechSmith.2023. -276c . ISBN 978-601-352-777-2.{[}in
Kazakh{]}

3. Fayaz M., Ahmad S., Ullah I., Kim D. A Blended Risk Index Modeling
and Visualization Based on Hierarchical Fuzzy Logic for Water Supply
Pipelines Assessment and Management//Processes. -2018. -Vol.6(5):61. DOI
10.3390/pr6050061.

4. Kafarov V.V., Dorohov I.N. Sistemnyj analiz processov himicheskoj
tehnologii. osnovy strategii: monografija//Izdatel' stvo
Jurajt. -2026.-499c. ISBN 978-5-534-06991-4. {[}in Russian{]}

5. Orazbayev B.B, Ospanov Ye.A., Orazbayeva K.N, Makhatova V.E.,
Urazgaliyeva M.K., Shagayeva A.B. Development of mathematical models of
R-1 reactor hydrotreatment unit using available information of various
types// \href{https://iopscience.iop.org/journal/1742-6596}{Journal of
Physics: Conference Series} 1399. -2019. -Vol.1399(4):044024. DOI
10.1088/1742-6596/1399/4/044024.

6. Grechuhina A.A., Elpidinskij A.A., Mingazov R.R., Plohova S.E. Raschet
rektifikacionnyh kolonn ustanovok peregonki nefti. Uchebnoe posobie
//Kazanskij nacional' nyj
issledovatel' skij tehnologicheskij universitet.
-2017.-92 c. ISBN 978-5-7882-2138-0. {[}in Russian{]}

7. Vlasov V.G. Podgotovka i pererabotka neftej. -2-e izd.//
Infra-Inzhenerija. -2026. -328 s. ISBN 978-5-9729-2808-8. {[}in
Russian{]}

8. Orazbayev B., Ospanov Ye., Makhatova V., Salybek L., Abdugulova Z,
Kulmagambetova Z., Suleimenova S., Orazbayeva K. Methods of Fuzzy
Multi-Criteria Decision Making for Controlling the Operating Modes of
the Stabilization Column of the Primary Oil-Refining
Unit//Mathematics.-2023.-Vol.11(13): 2820. DOI
\href{https://doi.org/10.3390/math11132820}{10.3390/math11132820}.

9. Muzipov H.N. Sistemy upravlenija tehnologicheskimi processami
dobychi, promyslovoj podgotovki i transporta nefti i gaza. Uchebnoe
posobie dlja vuzov // Izdatel' stvo
Lan'-Press. -2025. -268 c. ISBN 978-5-507-46261-2. {[}in
Russian{]}

10. Anwar M.Z., Al-Kenani A.N., Bashir S., Shabir M. Pessimistic
Multigranulation Rough Set of Intuitionistic Fuzzy Sets Based on Soft
Relations//Mathematics. -2022. -Vol.10(5):685. DOI 10.3390/
math10050685.

11. Merkulov V.V. Himicheskaja tehnologija. Pererabotka nefti i gaza:
Uchebnoe posobie// Izdatel' stvo Bastau. -2018. - 260 C.
ISBN 978-601-7275-69-6. {[}in Russian{]}

\emph{{\bfseries Сведения об авторах}}

Оразбаев Б.Б. - д.т.н., профессор, Евразийский национальный университет
им. Л.Н. Гумилева, Астана, Казахстан, e-mail:
batyr\_o@mail.ru;

Кужуханова Ж.А. - старший преподаватель, Казахский университет
технологии и бизнеса им. К.Кулажанова, Астана. Казахстан, e-mail:
kuzhuhanova-zhad@mail.ru;

Ускенбаева Г.А. - ассоциированный профессор, Евразийский национальный
университет им. Л.Н. Гумилева, Астана. Казахстан, e-mail:
gulzhum\_01@mail.ru;

Серимбетов Б.А. - ассоциированный профессор, Казахский университет
технологии и бизнеса им. К. Кулажанова, Астана. Казахстан, e-mail:
sba\_rnmc@mail.ru;

Курмангазиева Л.Т. - к.т.н., ассоциированный профессор, Атырауский
университет им. Х. Досмухамедова, Астана, Казахстан, e-mail:
kurmangazieval@mail.ru.

\emph{{\bfseries Information about the authors}}

Orazbaev B.B. - professor, L.N. Gumilyov Eurasian National University,
Astana, Kazakhstan, e-mail:
batyr\_o@mail.ru;

Kuzhuhanova Zh.A. - senior lecturer, K. Kulazhanov Kazakh University of
Technology and Business, Astana, Kazakhstan, e-mail:
\href{mailto:kuzhuhanova-zhad@mail.ru}{kuzhuhanoa-zhad@mail.ru};

Uskenbayeva G.A. - associate professor, L.N. Gumilyov Eurasian National
University, Astana, Kazakhstan, e-mail:
gulzhum\_01@mail.ru;

Serimbetov B.A. - associate professor, K. Kulazhanov Kazakh University
of Technology and Business, Astana, Kazakhstan, e-mail:
sba\_rnmc@mail.ru;

Kurmangaziyeva L.T. - associate professor, Kh.Dosmukhamedov Atyrau
University, Astana, Kazakhstan, e-mail:
kurmangazieval@mail.ru.\
